\documentclass[a3paper,portrait,11pt]{article}
\usepackage[margin=10mm]{geometry}
\usepackage[T1]{fontenc}
\usepackage[utf8]{inputenc}
\usepackage[french]{babel}
\usepackage{lmodern}
\usepackage{microtype}
\usepackage{amsmath,amssymb}
\usepackage[most]{tcolorbox}
\usepackage{tikz}
\usepackage{enumitem}
\usepackage{ragged2e}
\usetikzlibrary{arrows.meta,positioning,calc,shapes.geometric,shadows.blur,backgrounds,fit}

\pagestyle{empty}
\setlength{\parindent}{0pt}
\setlist[itemize]{leftmargin=5mm,itemsep=1.2mm,topsep=1mm}
\renewcommand{\familydefault}{\sfdefault}

% ----------------------------------------------------------------------------
% Palette
% ----------------------------------------------------------------------------
\definecolor{SAblue}{HTML}{0B3B78}
\definecolor{SAblue2}{HTML}{1769AA}
\definecolor{SAsky}{HTML}{EAF4FB}
\definecolor{SAgreen}{HTML}{5A9D4B}
\definecolor{SAgreenlight}{HTML}{EEF7EA}
\definecolor{SAorange}{HTML}{D88A22}
\definecolor{SAorangelight}{HTML}{FFF4E3}
\definecolor{SApurple}{HTML}{7251A8}
\definecolor{SApurplelight}{HTML}{F3EEFA}
\definecolor{SAred}{HTML}{C94747}
\definecolor{SAgray}{HTML}{4B5563}
\definecolor{SAlightgray}{HTML}{F5F7FA}

\newtcolorbox{SAbox}[2][]{
  enhanced,breakable,
  colback=white,colframe=#2,
  boxrule=1.1pt,arc=3mm,
  left=4mm,right=4mm,top=3.5mm,bottom=3.5mm,
  drop fuzzy shadow,
  #1
}

\newtcolorbox{SAtitlebox}[2][]{
  enhanced,
  colback=white,colframe=#2,
  colbacktitle=#2,coltitle=white,
  title={#1},fonttitle=\bfseries\Large,
  boxrule=1pt,arc=3mm,
  left=4mm,right=4mm,top=3mm,bottom=3mm,
  attach boxed title to top left={xshift=4mm,yshift*=-\tcboxedtitleheight/2},
  boxed title style={arc=2mm,boxrule=0pt,left=4mm,right=4mm},
  drop fuzzy shadow
}

\newcommand{\SApill}[2]{%
  \tikz[baseline=-0.55ex]\node[rounded corners=2mm,fill=#1!12,draw=#1,
  text=#1!80!black,inner xsep=3mm,inner ysep=1.2mm,font=\bfseries] {#2};%
}

\begin{document}

% ----------------------------------------------------------------------------
% Bandeau supérieur
% ----------------------------------------------------------------------------
\begin{tikzpicture}[remember picture,overlay]
  \fill[SAblue] (current page.north west) rectangle ([yshift=-44mm]current page.north east);
  \fill[SAblue2] ([xshift=-26mm,yshift=-44mm]current page.north east)
    circle[radius=42mm];
  \fill[white,opacity=.08] ([xshift=-5mm,yshift=-8mm]current page.north west)
    circle[radius=23mm];
  \draw[white,opacity=.22,line width=1.2pt]
    ([xshift=18mm,yshift=-13mm]current page.north west) circle[radius=7mm];
  \draw[white,opacity=.22,line width=1.2pt]
    ([xshift=29mm,yshift=-23mm]current page.north west) circle[radius=4mm];
\end{tikzpicture}

\vspace*{3mm}
\begin{minipage}[c]{0.74\textwidth}
  {\color{white}\fontsize{31}{34}\selectfont\bfseries Systèmes asservis}\par
  \vspace{2mm}
  {\color{white!88}\fontsize{15}{18}\selectfont
  Comprendre la boucle fermée, la modéliser et évaluer ses performances}
\end{minipage}\hfill
\begin{minipage}[c]{0.22\textwidth}
\begin{tikzpicture}[scale=.9,>=Latex]
  \draw[white,line width=1.8pt] (0,0) circle (1.55);
  \draw[white,line width=1.2pt] (0,0) circle (.85);
  \draw[white,line width=1.2pt] (-2,0)--(2,0);
  \draw[white,line width=1.2pt] (0,-2)--(0,2);
  \draw[white,line width=2pt,-{Stealth[length=4mm]}] (-1.55,-1.15)--(-.3,-.15);
  \fill[white] (0,0) circle (2.2pt);
\end{tikzpicture}
\end{minipage}

\vspace{13mm}

% ----------------------------------------------------------------------------
% Définition + relation fondamentale
% ----------------------------------------------------------------------------
\begin{minipage}[t]{0.58\textwidth}
\begin{SAtitlebox}[1. Définition]{SAblue2}
\justifying
Un \textbf{système asservi} ajuste automatiquement son action à partir d'une
\textbf{consigne} et d'une \textbf{mesure} de la sortie. Son objectif est de
réduire l'écart entre les deux, tout en restant performant malgré les
perturbations et les incertitudes du procédé.
\end{SAtitlebox}
\end{minipage}\hfill
\begin{minipage}[t]{0.39\textwidth}
\begin{SAbox}[colback=SAsky,halign=center,valign=center,height=47mm]{SAblue2}
{\color{SAblue}\Large\bfseries Relation fondamentale}\par\vspace{4mm}
{\fontsize{22}{25}\selectfont\bfseries
\(e(t)=r(t)-m(t)\)}\par\vspace{3mm}
{\color{SAgray}erreur = consigne $-$ mesure}
\end{SAbox}
\end{minipage}

\vspace{5mm}

% ----------------------------------------------------------------------------
% Schéma-bloc principal
% ----------------------------------------------------------------------------
\begin{SAtitlebox}[2. Structure d'une boucle fermée]{SAblue}
\begin{center}
\begin{tikzpicture}[
  >=Latex,
  block/.style={draw,rounded corners=2mm,minimum height=14mm,minimum width=35mm,
    align=center,font=\bfseries,blur shadow},
  signal/.style={-{Stealth[length=3.2mm]},line width=1.15pt,draw=SAgray},
  lab/.style={font=\bfseries\small,text=SAblue},
  sum/.style={circle,draw=SAblue,line width=1.2pt,minimum size=13mm,fill=SAsky}
]
  \node[lab] (r) at (-9.3,0) {Consigne\\$r(t)$};
  \node[sum] (s) at (-7.1,0) {};
  \node[block,fill=SAgreenlight,draw=SAgreen] (c) at (-3.8,0)
    {Correcteur\\$C(p)$};
  \node[block,fill=SAorangelight,draw=SAorange,minimum width=43mm] (g) at (1.0,0)
    {Actionneur + procédé\\$G(p)$};
  \node[lab] (y) at (6.0,0) {Sortie\\$y(t)$};
  \node[block,fill=SAsky,draw=SAblue2] (h) at (-1.1,-2.6)
    {Capteur\\$H(p)$};
  \node[lab,text=SAred] (d) at (1.0,2.0) {Perturbation\\$d(t)$};

  \draw[signal] (r.east)--(s.west);
  \draw[signal] (s.east)--node[above,font=\bfseries]{$e(t)$}(c.west);
  \draw[signal] (c.east)--node[above,font=\small]{commande}(g.west);
  \draw[signal] (g.east)--coordinate (pick)(y.west);
  \draw[signal,draw=SAred] (d.south)--(g.north);
  \draw[signal] (pick) |- (h.east);
  \draw[signal] (h.west) -| node[pos=.17,below,font=\small\bfseries,text=SAblue2]{mesure $m(t)$} (s.south);

  \node at ($(s.center)+(-.27,.28)$) {\bfseries +};
  \node at ($(s.center)+(-.25,-.30)$) {\bfseries --};

  \node[draw=SAblue2,dashed,rounded corners=2mm,fit=(s)(c)(g)(h),
    inner xsep=5mm,inner ysep=4mm,label={[text=SAblue2,font=\bfseries]below:la mesure referme la boucle}] {};
\end{tikzpicture}
\end{center}

\vspace{1mm}
\begin{center}
\SApill{SAblue2}{Comparer}\quad
\SApill{SAgreen}{Corriger}\quad
\SApill{SAorange}{Agir}\quad
\SApill{SAblue}{Mesurer}\quad
\SApill{SAred}{Rejeter les perturbations}
\end{center}
\end{SAtitlebox}

\vspace{5mm}

% ----------------------------------------------------------------------------
% Performances
% ----------------------------------------------------------------------------
\begin{center}
{\color{SAblue}\LARGE\bfseries 3. Les trois performances fondamentales}
\end{center}
\vspace{2mm}

\begin{minipage}[t]{0.32\textwidth}
\begin{SAbox}[height=59mm,colback=SAgreenlight]{SAgreen}
\begin{center}
\begin{tikzpicture}
  \draw[SAgreen,line width=1.6pt,rounded corners=1mm] (-.8,-.9)--(-.8,.75)--(0,1.1)--(.8,.75)--(.8,-.9);
  \draw[SAgreen,line width=1.6pt] (-.5,.05)--(-.1,-.35)--(.55,.45);
\end{tikzpicture}\par
{\color{SAgreen!70!black}\Large\bfseries Stabilité}\par\vspace{2mm}
Le système reste borné et rejoint un régime permanent après une sollicitation.
\end{center}
\end{SAbox}
\end{minipage}\hfill
\begin{minipage}[t]{0.32\textwidth}
\begin{SAbox}[height=59mm,colback=SAsky]{SAblue2}
\begin{center}
\begin{tikzpicture}
  \draw[SAblue2,line width=1.6pt] (0,0) circle (1);
  \draw[SAblue2,line width=1.6pt] (0,0) circle (.58);
  \draw[SAblue2,line width=1.6pt] (0,0) circle (.18);
  \draw[SAblue2,line width=1.6pt,-{Stealth[length=3mm]}] (-1.2,-1.0)--(-.1,-.1);
\end{tikzpicture}\par
{\color{SAblue2}\Large\bfseries Précision}\par\vspace{2mm}
La sortie atteint la consigne avec une erreur statique aussi faible que possible.
\end{center}
\end{SAbox}
\end{minipage}\hfill
\begin{minipage}[t]{0.32\textwidth}
\begin{SAbox}[height=59mm,colback=SApurplelight]{SApurple}
\begin{center}
\begin{tikzpicture}
  \draw[SApurple,line width=1.6pt] (0,0) circle (.95);
  \draw[SApurple,line width=1.6pt] (0,.95)--(0,1.25);
  \draw[SApurple,line width=1.6pt] (-.25,1.25)--(.25,1.25);
  \draw[SApurple,line width=1.6pt,-{Stealth[length=2.5mm]}] (0,0)--(.45,.48);
  \draw[SApurple,line width=1.4pt] (0,0)--(0,-.55);
\end{tikzpicture}\par
{\color{SApurple}\Large\bfseries Rapidité}\par\vspace{2mm}
Le temps de montée et le temps de réponse restent compatibles avec le besoin.
\end{center}
\end{SAbox}
\end{minipage}

\vspace{5mm}

% ----------------------------------------------------------------------------
% Modèles + démarche
% ----------------------------------------------------------------------------
\begin{minipage}[t]{0.49\textwidth}
\begin{SAtitlebox}[4. Modèles essentiels]{SApurple}
\begin{tabular}{@{}p{.34\linewidth}p{.60\linewidth}@{}}
\textbf{Fonction de transfert} & $H(p)=\dfrac{S(p)}{E(p)}$\\[3mm]
\textbf{Boucle fermée} & $H_{BF}(p)=\dfrac{G(p)}{1+G(p)H(p)}$\\[3mm]
\textbf{Premier ordre} & $H(p)=\dfrac{K}{1+\tau p}$\\[3mm]
\textbf{Second ordre} & $H(p)=\dfrac{K\omega_0^2}{p^2+2\xi\omega_0p+\omega_0^2}$
\end{tabular}
\end{SAtitlebox}
\end{minipage}\hfill
\begin{minipage}[t]{0.49\textwidth}
\begin{SAtitlebox}[5. Démarche d'étude]{SAorange}
\begin{enumerate}[leftmargin=8mm,itemsep=1.4mm,topsep=0mm]
  \item Identifier consigne, sortie, perturbations et capteur.
  \item Construire le schéma-bloc et déterminer la FTBO.
  \item Calculer la FTBF et repérer pôles et zéros.
  \item Étudier les réponses temporelle et fréquentielle.
  \item Conclure sur stabilité, précision, rapidité et robustesse.
\end{enumerate}
\end{SAtitlebox}
\end{minipage}

\vspace{5mm}

% ----------------------------------------------------------------------------
% Exemples
% ----------------------------------------------------------------------------
\begin{SAtitlebox}[6. Exemples de systèmes asservis]{SAgreen}
\begin{center}
\begin{tikzpicture}[
  ex/.style={rounded corners=3mm,draw=SAgreen,fill=SAgreenlight,
    minimum width=70mm,minimum height=22mm,align=center,font=\bfseries},
  arr/.style={-{Stealth[length=3mm]},line width=1.2pt,draw=SAgreen}
]
  \node[ex] (a) at (-8,0) {Régulateur de vitesse\\\small maintenir une vitesse malgré la pente};
  \node[ex] (b) at (0,0) {Thermostat\\\small maintenir une température malgré les pertes};
  \node[ex] (c) at (8,0) {Axe robotisé\\\small atteindre une position malgré la charge};
  \draw[arr] ($(a.south)+(-1,-.35)$)--($(a.south)+(1,-.35)$);
  \draw[arr] ($(b.south)+(-1,-.35)$)--($(b.south)+(1,-.35)$);
  \draw[arr] ($(c.south)+(-1,-.35)$)--($(c.south)+(1,-.35)$);
\end{tikzpicture}
\end{center}
\end{SAtitlebox}

\vfill
\begin{tcolorbox}[
  enhanced,colback=SAblue,colframe=SAblue,arc=3mm,boxrule=0pt,
  left=6mm,right=6mm,top=3.5mm,bottom=3.5mm]
\begin{center}
{\color{white}\Large\bfseries À retenir : atteindre la consigne avec le meilleur compromis}
\quad
{\color{white!90}\large stabilité -- précision -- rapidité}
\end{center}
\end{tcolorbox}

\end{document}
